\documentclass[12pt,a4paper]{article}
\usepackage[utf8]{inputenc}
\usepackage[T1]{fontenc}
\usepackage[spanish,es-noshorthands]{babel}
\usepackage{geometry}
\usepackage{graphicx}
\usepackage{xcolor}
\usepackage{listings}
\usepackage{hyperref}
\usepackage{tabularx}
\usepackage{booktabs}
\usepackage{enumitem}
\usepackage{tikz}
\usepackage{fancyhdr}
\usepackage{titlesec}
\usepackage{multicol}

\geometry{margin=2.5cm}

% Colores del proyecto
\definecolor{darkblue}{HTML}{001F3F}
\definecolor{primarygreen}{HTML}{00DF5F}
\definecolor{codebg}{HTML}{F3F3F4}
\definecolor{codetext}{HTML}{2D3748}

\hypersetup{
    colorlinks=true,
    linkcolor=darkblue,
    urlcolor=primarygreen,
}

% Estilo de codigo
\lstset{
    backgroundcolor=\color{codebg},
    basicstyle=\ttfamily\small\color{codetext},
    breaklines=true,
    frame=single,
    rulecolor=\color{gray!30},
    numbers=left,
    numberstyle=\tiny\color{gray},
    tabsize=2,
}

% Header/Footer
\pagestyle{fancy}
\fancyhf{}
\fancyhead[L]{\textcolor{darkblue}{\textbf{CONOCEMEX}}}
\fancyhead[R]{\textcolor{gray}{Arquitectura de Software}}
\fancyfoot[C]{\thepage}

\titleformat{\section}{\Large\bfseries\color{darkblue}}{}{0em}{}
\titleformat{\subsection}{\large\bfseries\color{darkblue!80}}{}{0em}{}
\titleformat{\subsubsection}{\normalsize\bfseries\color{darkblue!60}}{}{0em}{}

\begin{document}

% ================================================================
% PORTADA
% ================================================================
\begin{titlepage}
\centering
\vspace*{3cm}

{\Huge\bfseries\textcolor{darkblue}{CONOCEMEX}\par}
\vspace{0.5cm}
{\Large\textcolor{primarygreen}{Arquitectura de Software}\par}
\vspace{1cm}
{\large App M\'ovil del Microempresario (Flutter)\par}

\vspace{2cm}

\begin{tabular}{ll}
\textbf{Equipo:} & JaguarMind \\
\textbf{Evento:} & Hackathon Talent Land M\'exico 2026 \\
\textbf{Stack:} & Flutter + Supabase + Google Cloud \\
\textbf{Versi\'on:} & 1.0 \\
\textbf{Fecha:} & Abril 2026 \\
\end{tabular}

\vfill

{\small Documento generado autom\'aticamente a partir del c\'odigo fuente.}
\end{titlepage}

% ================================================================
% INDICE
% ================================================================
\tableofcontents
\newpage

% ================================================================
% 1. VISION GENERAL
% ================================================================
\section{Visi\'on General del Proyecto}

CONOCEMEX es una plataforma tur\'istica que conecta \textbf{microempresarios mexicanos} con \textbf{turistas internacionales}. La aplicaci\'on m\'ovil Flutter est\'a dise\~nada para los vendedores/microempresarios, permiti\'endoles:

\begin{itemize}[leftmargin=2cm]
    \item Registrar y gestionar sus negocios
    \item Agregar productos y servicios con fotos e IA
    \item Recibir pagos v\'ia QR de Mercado Pago
    \item Comunicarse con turistas en su idioma (traducci\'on autom\'atica)
    \item Participar en comunidades tem\'aticas
\end{itemize}

\subsection{Datos T\'ecnicos}

\begin{tabularx}{\textwidth}{lX}
\toprule
\textbf{M\'etrica} & \textbf{Valor} \\
\midrule
Archivos Dart & 89 \\
M\'odulos de features & 6 (Auth, Profile, Business, Category, Offering, Chat) \\
Servicios registrados & 14 \\
Rutas de navegaci\'on & 9 \\
Use Cases & 13 \\
ViewModels & 4 \\
Idiomas soportados & 4 (ES, EN, FR, PT) \\
APIs externas integradas & 7 \\
\bottomrule
\end{tabularx}

% ================================================================
% 2. ARQUITECTURA
% ================================================================
\section{Arquitectura de Software}

\subsection{Patr\'on: Clean Architecture}

La aplicaci\'on sigue el patr\'on de \textbf{Clean Architecture} con tres capas claramente separadas:

\begin{center}
\begin{tikzpicture}[
    layer/.style={draw=darkblue, fill=darkblue!5, rounded corners=8pt, minimum width=12cm, minimum height=1.8cm, align=center, font=\bfseries},
    arrow/.style={->, thick, darkblue!60}
]
    \node[layer, fill=primarygreen!10] (pres) at (0, 4) {\textcolor{darkblue}{Capa de Presentaci\'on}\\[-2pt]\small\normalfont Pages, Widgets, ViewModels};
    \node[layer, fill=darkblue!8] (domain) at (0, 2) {\textcolor{darkblue}{Capa de Dominio}\\[-2pt]\small\normalfont Entities, Repositories (abstract), Use Cases};
    \node[layer, fill=codebg] (data) at (0, 0) {\textcolor{darkblue}{Capa de Datos}\\[-2pt]\small\normalfont Models, DataSources, Repository Implementations};

    \draw[arrow] (pres.south) -- (domain.north) node[midway, right, font=\small] {depende de};
    \draw[arrow] (data.north) -- (domain.south) node[midway, right, font=\small] {implementa};
\end{tikzpicture}
\end{center}

\subsection{Gesti\'on de Estado: Provider + ChangeNotifier}

\begin{tabularx}{\textwidth}{lX}
\toprule
\textbf{ViewModel} & \textbf{Responsabilidad} \\
\midrule
LoginViewModel & Autenticaci\'on (login, signup, Google, biom\'etrico, logout) \\
DashboardViewModel & Carga de perfil, lista de negocios, actualizaci\'on de rol \\
CreateBusinessViewModel & Creaci\'on de negocio, carga de categor\'ias \\
OfferingsViewModel & CRUD de productos/servicios por negocio \\
\bottomrule
\end{tabularx}

\subsection{Inyecci\'on de Dependencias: GetIt}

Todos los servicios se registran como \textbf{singletons} en \texttt{setup\_dependencies.dart}. El contenedor se inicializa antes del primer frame de la UI.

% ================================================================
% 3. ESTRUCTURA DE ARCHIVOS
% ================================================================
\section{Estructura de Archivos}

\begin{lstlisting}[language=bash, title=Estructura del proyecto]
lib/
|-- core/
|   |-- config/env.dart              # Variables de entorno
|   |-- constants/app_constants.dart # Rutas y constantes
|   |-- di/setup_dependencies.dart   # Registro de DI
|   |-- extensions/exceptions.dart   # Manejo de errores
|   |-- routes/app_routes.dart       # Generador de rutas
|   |-- services/
|       |-- biometric_service.dart
|       |-- cloudinary_service.dart
|       |-- community_chat_service.dart
|       |-- community_service.dart
|       |-- deepl_service.dart
|       |-- gemini_service.dart
|       |-- image_picker_service.dart
|       |-- locale_service.dart
|       |-- mercado_pago_service.dart
|       |-- mp_charge_service.dart
|       |-- secure_storage_service.dart
|-- features/
|   |-- auth/        # Autenticacion y navegacion principal
|   |-- profile/     # Gestion de perfil de usuario
|   |-- business/    # Gestion de negocios
|   |-- category/    # Catalogo de categorias
|   |-- offering/    # Productos y servicios
|-- l10n/            # Archivos de internacionalizacion
|   |-- app_es.arb, app_en.arb, app_fr.arb, app_pt.arb
|-- main.dart        # Punto de entrada
\end{lstlisting}

% ================================================================
% 4. MODULOS
% ================================================================
\section{M\'odulos del Sistema}

\subsection{M\'odulo de Autenticaci\'on}

Soporta tres m\'etodos de autenticaci\'on:

\begin{enumerate}
    \item \textbf{Google Sign-In}: OAuth 2.0 con \texttt{google\_sign\_in} + \texttt{signInWithIdToken} de Supabase
    \item \textbf{Email/Password}: \texttt{signInWithPassword} y \texttt{signUp} de Supabase Auth
    \item \textbf{Biom\'etrico}: \texttt{local\_auth} para huella dactilar / reconocimiento facial
\end{enumerate}

\begin{center}
\begin{tikzpicture}[
    box/.style={draw=darkblue, fill=codebg, rounded corners=4pt, minimum width=3cm, minimum height=1cm, align=center, font=\small},
    arrow/.style={->, thick, primarygreen!80!black}
]
    \node[box] (login) at (0, 0) {LoginPage};
    \node[box] (signup) at (4, 0) {SignUpPage};
    \node[box] (onboard) at (8, 0) {OnboardingPage};
    \node[box] (home) at (12, 0) {MainShellPage};

    \draw[arrow] (login) -- (signup) node[midway, above, font=\tiny] {Sign Up};
    \draw[arrow] (login) -- (onboard) node[midway, below, font=\tiny] {Login OK};
    \draw[arrow] (signup) -- (onboard) node[midway, above, font=\tiny] {Registro OK};
    \draw[arrow] (onboard) -- (home) node[midway, above, font=\tiny] {Perfil completo};
\end{tikzpicture}
\end{center}

\subsection{M\'odulo de Negocios}

\begin{tabularx}{\textwidth}{lX}
\toprule
\textbf{Funcionalidad} & \textbf{Descripci\'on} \\
\midrule
Registro de negocio & Formulario con foto (Cloudinary), IA (OpenRouter), mapa (Google Maps) \\
Selecci\'on en mapa & Google Maps con reverse geocoding para obtener direcci\'on autom\'atica \\
Autocompletado IA & An\'alisis de imagen del local para extraer nombre, categor\'ia, direcci\'on \\
Gesti\'on de productos & CRUD de offerings con foto, precio, tipo, duraci\'on \\
\bottomrule
\end{tabularx}

\subsection{M\'odulo de Pagos (Mercado Pago)}

\begin{center}
\begin{tikzpicture}[
    box/.style={draw=darkblue, fill=codebg, rounded corners=4pt, minimum width=3.5cm, minimum height=1.2cm, align=center, font=\small},
    arrow/.style={->, thick, darkblue!60}
]
    \node[box] (connect) at (0, 0) {1. Conectar MP\\(OAuth)};
    \node[box] (charge) at (5, 0) {2. Generar cobro\\(Edge Function)};
    \node[box] (qr) at (10, 0) {3. Mostrar QR\\(qr\_flutter)};
    \node[box] (pay) at (5, -2.5) {4. Turista paga\\(Checkout Pro)};
    \node[box] (done) at (10, -2.5) {5. Pago recibido\\(Realtime)};

    \draw[arrow] (connect) -- (charge);
    \draw[arrow] (charge) -- (qr);
    \draw[arrow] (qr) -- (pay);
    \draw[arrow] (pay) -- (done);
\end{tikzpicture}
\end{center}

\subsection{M\'odulo de Chat}

El sistema de chat tiene dos modalidades:

\subsubsection{Chat Grupal (Comunidades)}
\begin{itemize}
    \item Crear comunidades con c\'odigo de invitaci\'on
    \item Mensajes en tiempo real v\'ia Supabase Realtime
    \item Traducci\'on autom\'atica al idioma del usuario (DeepL API)
    \item Cache de traducciones en memoria
\end{itemize}

\subsubsection{Chat 1 a 1 (Turista $\leftrightarrow$ Vendedor)}
\begin{itemize}
    \item El turista inicia la conversaci\'on desde la ficha del negocio
    \item El vendedor ve las conversaciones en su lista de chats
    \item Contador de mensajes no le\'idos (badges)
    \item Traducci\'on autom\'atica bidireccional
\end{itemize}

% ================================================================
% 5. INTEGRACIONES EXTERNAS
% ================================================================
\section{Integraciones Externas}

\begin{tabularx}{\textwidth}{llX}
\toprule
\textbf{Servicio} & \textbf{Prop\'osito} & \textbf{Implementaci\'on} \\
\midrule
Supabase & Backend completo & Auth, PostgreSQL, Realtime, Edge Functions \\
Google Sign-In & Autenticaci\'on OAuth & \texttt{google\_sign\_in} + Supabase \texttt{signInWithIdToken} \\
Google Maps & Selecci\'on de ubicaci\'on & \texttt{google\_maps\_flutter} + Geocoding API \\
Cloudinary & Almacenamiento de im\'agenes & Upload unsigned preset, transformaciones CDN \\
OpenRouter & IA visi\'on (an\'alisis de fotos) & Gemini 2.0 Flash v\'ia API compatible OpenAI \\
DeepL & Traducci\'on autom\'atica & API Free, cache en memoria, 4 idiomas \\
Mercado Pago & Pagos QR & OAuth + Checkout Pro + Webhook \\
\bottomrule
\end{tabularx}

% ================================================================
% 6. INTERNACIONALIZACION
% ================================================================
\section{Internacionalizaci\'on (i18n)}

\subsection{Idiomas Soportados}

\begin{tabularx}{\textwidth}{llll}
\toprule
\textbf{C\'odigo} & \textbf{Idioma} & \textbf{Archivo ARB} & \textbf{C\'odigo DeepL} \\
\midrule
es & Espa\~nol & app\_es.arb & ES \\
en & English & app\_en.arb & EN \\
fr & Fran\c{c}ais & app\_fr.arb & FR \\
pt & Portugu\^es & app\_pt.arb & PT-BR \\
\bottomrule
\end{tabularx}

\subsection{Flujo de Cambio de Idioma}

\begin{enumerate}
    \item Usuario selecciona idioma en onboarding o edici\'on de perfil
    \item Se guarda en \texttt{profile\_languages} (Supabase) con \texttt{is\_preferred = true}
    \item \texttt{LocaleService} actualiza el locale activo y lo persiste en \texttt{SharedPreferences}
    \item \texttt{MaterialApp} reconstruye la UI completa en el nuevo idioma (sin reiniciar)
\end{enumerate}

% ================================================================
% 7. BASE DE DATOS
% ================================================================
\section{Base de Datos (Supabase / PostgreSQL)}

\subsection{Tablas Principales}

\begin{tabularx}{\textwidth}{llX}
\toprule
\textbf{Tabla} & \textbf{M\'odulo} & \textbf{Descripci\'on} \\
\midrule
profiles & Auth/Profile & Extensi\'on de auth.users con rol, nombre, nacionalidad \\
businesses & Business & Perfil digital del micronegocio con PostGIS \\
categories & Category & 8 giros tur\'isticos con traducciones \\
offerings & Offering & Productos y servicios con precio y tipo \\
offering\_translations & Offering & Nombre/descripci\'on multiidioma \\
communities & Chat & Grupos de chat con c\'odigo de invitaci\'on \\
community\_members & Chat & Membres\'ia N:M con roles \\
community\_messages & Chat & Mensajes en tiempo real \\
conversations & Chat 1:1 & Conversaciones directas turista-vendedor \\
direct\_messages & Chat 1:1 & Mensajes del chat directo \\
transactions & Pagos & Cobros con estado y referencia a MP \\
mp\_accounts & Pagos & Tokens OAuth de Mercado Pago por negocio \\
\bottomrule
\end{tabularx}

\subsection{Algoritmo de Recomendaci\'on}

La base de datos incluye dos funciones RPC para el feed del turista:

\begin{itemize}
    \item \texttt{recommend\_businesses\_discovery}: Ranking personalizado con 3 slots de equidad (reduce Gini de 0.52 a 0.33)
    \item \texttt{recommend\_businesses\_top\_quality}: Ranking puro por calidad sin boost de equidad
\end{itemize}

% ================================================================
% 8. SEGURIDAD
% ================================================================
\section{Seguridad}

\begin{tabularx}{\textwidth}{lX}
\toprule
\textbf{Mecanismo} & \textbf{Implementaci\'on} \\
\midrule
Autenticaci\'on & Supabase Auth con JWT, Google OAuth 2.0, biom\'etrico \\
Autorizaci\'on & Row Level Security (RLS) en todas las tablas \\
Tokens de pago & Cifrados en \texttt{mp\_accounts}, nunca expuestos al cliente \\
Upload de im\'agenes & Preset unsigned de Cloudinary (solo crear, no borrar) \\
Variables sensibles & Archivo \texttt{.env} excluido de control de versiones \\
Almacenamiento local & \texttt{SharedPreferences} para tokens de sesi\'on \\
\bottomrule
\end{tabularx}

% ================================================================
% 9. DEPENDENCIAS
% ================================================================
\section{Dependencias del Proyecto}

\subsection{Dependencias de Producci\'on}

\begin{multicols}{2}
\begin{itemize}[leftmargin=0.5cm, itemsep=1pt]
    \item \texttt{supabase\_flutter} --- Backend
    \item \texttt{google\_sign\_in} --- Auth Google
    \item \texttt{local\_auth} --- Biom\'etrico
    \item \texttt{provider} --- State management
    \item \texttt{get\_it} --- Inyecci\'on de dependencias
    \item \texttt{dio} --- Cliente HTTP
    \item \texttt{google\_maps\_flutter} --- Mapas
    \item \texttt{qr\_flutter} --- C\'odigos QR
    \item \texttt{image\_picker} --- Selecci\'on de fotos
    \item \texttt{url\_launcher} --- Abrir URLs
    \item \texttt{share\_plus} --- Compartir nativo
    \item \texttt{shared\_preferences} --- Storage local
    \item \texttt{flutter\_dotenv} --- Variables .env
    \item \texttt{connectivity\_plus} --- Red
    \item \texttt{equatable} --- Igualdad de objetos
    \item \texttt{intl} --- Internacionalizaci\'on
    \item \texttt{flutter\_localizations} --- i18n
\end{itemize}
\end{multicols}

% ================================================================
% 10. FLUJO DE USUARIO
% ================================================================
\section{Flujo Principal del Usuario}

\begin{center}
\begin{tikzpicture}[
    step/.style={draw=darkblue, fill=primarygreen!10, rounded corners=6pt, minimum width=4cm, minimum height=1cm, align=center, font=\small\bfseries},
    arrow/.style={->, thick, darkblue!60},
    note/.style={font=\tiny\itshape, text=gray}
]
    \node[step] (splash) at (0, 0) {Splash};
    \node[step] (login) at (0, -1.8) {Login / SignUp};
    \node[step] (onboard) at (0, -3.6) {Onboarding};
    \node[step] (home) at (0, -5.4) {Mi Panel (Home)};

    \node[step] (create) at (-5, -7.2) {Crear Negocio};
    \node[step] (catalog) at (0, -7.2) {Cat\'alogo};
    \node[step] (chat) at (5, -7.2) {Chat};

    \node[step] (offering) at (-5, -9) {Agregar Producto};
    \node[step] (cobrar) at (0, -9) {Cobrar (QR)};
    \node[step] (community) at (5, -9) {Comunidades};

    \draw[arrow] (splash) -- (login) node[midway, right, note] {sin sesi\'on};
    \draw[arrow] (login) -- (onboard) node[midway, right, note] {auth OK};
    \draw[arrow] (onboard) -- (home) node[midway, right, note] {perfil completo};
    \draw[arrow] (home) -- (create);
    \draw[arrow] (home) -- (catalog);
    \draw[arrow] (home) -- (chat);
    \draw[arrow] (create) -- (offering);
    \draw[arrow] (catalog) -- (cobrar);
    \draw[arrow] (chat) -- (community);
\end{tikzpicture}
\end{center}

% ================================================================
% 11. NAVEGACION
% ================================================================
\section{Navegaci\'on}

\subsection{Bottom Navigation Bar (MainShellPage)}

La aplicaci\'on utiliza un \texttt{IndexedStack} con 4 tabs:

\begin{tabularx}{\textwidth}{clX}
\toprule
\textbf{Tab} & \textbf{Icono} & \textbf{Contenido} \\
\midrule
0 & Home & Dashboard de negocios (Mi Panel) \\
1 & Grid & Cat\'alogo de productos + Pagos MP \\
2 & Chat & Chats directos + Comunidades \\
3 & Person & Perfil del usuario con edici\'on \\
\bottomrule
\end{tabularx}

\subsection{Rutas Registradas}

\begin{tabularx}{\textwidth}{lX}
\toprule
\textbf{Ruta} & \textbf{P\'agina} \\
\midrule
\texttt{/} & SplashPage (verificaci\'on de sesi\'on) \\
\texttt{/login} & LoginPage (Google + Email/Password) \\
\texttt{/signup} & SignUpPage (registro de vendedor) \\
\texttt{/onboarding} & OnboardingPage (completar perfil) \\
\texttt{/home} & MainShellPage (4 tabs) \\
\texttt{/profile/edit} & EditProfilePage \\
\texttt{/business/create} & CreateBusinessPage (con mapa e IA) \\
\texttt{/business/detail} & BusinessDetailPage \\
\texttt{/offering/create} & CreateOfferingPage \\
\bottomrule
\end{tabularx}

% ================================================================
% 12. DISENO
% ================================================================
\section{Dise\~no Visual (Sistema Stitch)}

\subsection{Paleta de Colores}

\begin{tabularx}{\textwidth}{llX}
\toprule
\textbf{Variable} & \textbf{Hex} & \textbf{Uso} \\
\midrule
\texttt{darkBlue} & \texttt{\#001F3F} & Textos principales, t\'itulos, iconos \\
\texttt{primaryGreen} & \texttt{\#00DF5F} & Botones primarios, acentos, badges activos \\
\texttt{bgGrey} & \texttt{\#F3F3F4} & Fondos de inputs, burbujas de chat ajenas \\
\texttt{white} & \texttt{\#FFFFFF} & Fondo de pantallas y cards \\
\bottomrule
\end{tabularx}

\subsection{Componentes UI}

\begin{itemize}
    \item \textbf{Botones primarios}: Pill shape (\texttt{borderRadius: 28}), verde con sombra
    \item \textbf{Inputs}: Fondo gris sin borde, focus verde, labels uppercase con letter-spacing
    \item \textbf{Cards}: Sin elevaci\'on, borde sutil (\texttt{alpha: 0.06})
    \item \textbf{Bottom sheets}: Bordes redondeados superiores con handle visual
    \item \textbf{Tipograf\'ia}: Weight 900 para t\'itulos, 700 para labels, 500/600 para body
\end{itemize}

% ================================================================
% CONCLUSION
% ================================================================
\section{Conclusi\'on}

CONOCEMEX representa una soluci\'on integral para la digitalizaci\'on de micronegocios tur\'isticos en M\'exico. La arquitectura limpia, la integraci\'on de m\'ultiples servicios en la nube, y el soporte multiidioma con traducci\'on autom\'atica permiten que vendedores de cualquier parte del pa\'is puedan ofrecer sus productos y servicios a turistas de todo el mundo, eliminando barreras de idioma y facilitando los pagos digitales.

\vspace{1cm}

\begin{center}
\textcolor{darkblue}{\textbf{Equipo JaguarMind --- Talent Land M\'exico 2026}}
\end{center}

\end{document}
